%%%%%%%%%%%%%%%%%%%%%%%%%%%%%%%%%%%%%%%%%%%%%%%%%%%%%%%%%%%%%%%%%%%%%%%%%%%%%%%%
% ACF-DOC-039.04
% Heatwave and Drought Monitoring System
%%%%%%%%%%%%%%%%%%%%%%%%%%%%%%%%%%%%%%%%%%%%%%%%%%%%%%%%%%%%%%%%%%%%%%%%%%%%%%%%

\section{Heatwave and Drought Monitoring System}

\subsection{Executive Summary}

The ACF Heatwave and Drought Monitoring System is a global climate
intelligence module dedicated to the detection, prediction and analysis
of extreme temperature events and water scarcity conditions.

Heatwaves and droughts are complex hazards resulting from interactions
between:

\begin{itemize}

\item Atmospheric circulation

\item Surface energy balance

\item Soil moisture conditions

\item Ocean-atmosphere interactions

\item Climate variability

\item Human activities

\end{itemize}


The system integrates:

\begin{itemize}

\item Numerical Weather Prediction

\item Climate models

\item Satellite observations

\item Land surface models

\item Soil moisture analysis

\item Artificial intelligence

\item Earth System Digital Twin

\end{itemize}


The objectives are:

\begin{itemize}

\item Early detection of heatwave events

\item Drought monitoring

\item Agricultural risk assessment

\item Water resource management

\item Emergency decision support

\end{itemize}


%%%%%%%%%%%%%%%%%%%%%%%%%%%%%%%%%%%%%%%%%%%%%%%%%%%%%%%%%%%%%%%%%%%%%%%%%%%%%%%%

\subsection{Introduction}

Extreme heat and drought represent major global environmental hazards.

Heatwaves are periods of abnormally high temperatures lasting several
days or longer.

Droughts correspond to prolonged periods with insufficient water
availability compared with normal conditions.


The ACF approach considers these hazards as coupled Earth system
processes.


The hazard chain is:

\[
Climate
\rightarrow
Atmosphere
\rightarrow
Land Surface
\rightarrow
Water Resources
\rightarrow
Impact
\]


%%%%%%%%%%%%%%%%%%%%%%%%%%%%%%%%%%%%%%%%%%%%%%%%%%%%%%%%%%%%%%%%%%%%%%%%%%%%%%%%

\subsection{Heatwave Detection Concept}

The ACF Heatwave Engine detects extreme temperature events using:

\begin{itemize}

\item Temperature anomalies

\item Persistence duration

\item Geographic extension

\item Human thermal stress

\item Historical climatology

\end{itemize}


A heatwave anomaly is calculated as:

\[
T_{anom}
=
T_{observed}
-
T_{climatology}
\]


where:

\begin{itemize}

\item $T_{observed}$ is current temperature

\item $T_{climatology}$ is long-term reference temperature

\end{itemize}


%%%%%%%%%%%%%%%%%%%%%%%%%%%%%%%%%%%%%%%%%%%%%%%%%%%%%%%%%%%%%%%%%%%%%%%%%%%%%%%%

\subsection{Temperature Monitoring}

The system analyzes:

\begin{itemize}

\item Surface air temperature

\item Maximum temperature

\item Minimum temperature

\item Night-time temperature

\item Thermal persistence

\end{itemize}


The ACF temperature intelligence combines:

\begin{itemize}

\item Weather stations

\item Satellite land surface temperature

\item Numerical models

\item Reanalysis datasets

\end{itemize}


%%%%%%%%%%%%%%%%%%%%%%%%%%%%%%%%%%%%%%%%%%%%%%%%%%%%%%%%%%%%%%%%%%%%%%%%%%%%%%%%
%%%%%%%%%%%%%%%%%%%%%%%%%%%%%%%%%%%%%%%%%%%%%%%%%%%%%%%%%%%%%%%%%%%%%%%%%%%%%%%%
\subsection{Heat Stress Indices}

Temperature alone does not fully describe human thermal stress.

The ACF Heatwave Engine calculates multiple thermal comfort indices.

The main indicators are:

\begin{itemize}

\item Heat Index (HI)

\item Wet Bulb Globe Temperature (WBGT)

\item Universal Thermal Climate Index (UTCI)

\item Apparent Temperature

\item Humidex

\end{itemize}


%%%%%%%%%%%%%%%%%%%%%%%%%%%%%%%%%%%%%%%%%%%%%%%%%%%%%%%%%%%%%%%%%%%%%%%%%%%%%%%%

\subsubsection{Heat Index}

The Heat Index combines air temperature and humidity.

A simplified formulation is:

\[
HI=f(T,RH)
\]

where:

\begin{itemize}

\item $T$ = air temperature

\item $RH$ = relative humidity

\end{itemize}


High Heat Index values indicate increased physiological stress.


%%%%%%%%%%%%%%%%%%%%%%%%%%%%%%%%%%%%%%%%%%%%%%%%%%%%%%%%%%%%%%%%%%%%%%%%%%%%%%%%

\subsubsection{Wet Bulb Globe Temperature}

WBGT evaluates thermal stress by combining:

\begin{itemize}

\item Air temperature

\item Humidity

\item Radiation

\item Wind conditions

\end{itemize}


The simplified equation is:

\[
WBGT=
0.7T_w+
0.2T_g+
0.1T_a
\]


where:

\begin{itemize}

\item $T_w$ = wet bulb temperature

\item $T_g$ = globe temperature

\item $T_a$ = air temperature

\end{itemize}


%%%%%%%%%%%%%%%%%%%%%%%%%%%%%%%%%%%%%%%%%%%%%%%%%%%%%%%%%%%%%%%%%%%%%%%%%%%%%%%%

\subsubsection{Universal Thermal Climate Index}

UTCI estimates human thermal response under outdoor conditions.

It includes:

\begin{itemize}

\item Temperature

\item Wind speed

\item Humidity

\item Radiation

\end{itemize}


The ACF system generates global thermal stress maps.

%%%%%%%%%%%%%%%%%%%%%%%%%%%%%%%%%%%%%%%%%%%%%%%%%%%%%%%%%%%%%%%%%%%%%%%%%%%%%%%%

\subsection{Surface Energy Balance}

Heatwaves are strongly controlled by the surface energy balance.

The fundamental equation is:

\[
R_n=
H+
LE+
G
\]


where:

\begin{itemize}

\item $R_n$ = net radiation

\item $H$ = sensible heat flux

\item $LE$ = latent heat flux

\item $G$ = ground heat flux

\end{itemize}


The ACF land surface module evaluates:

\begin{itemize}

\item Soil moisture

\item Vegetation state

\item Evapotranspiration

\item Surface temperature

\item Radiation balance

\end{itemize}


%%%%%%%%%%%%%%%%%%%%%%%%%%%%%%%%%%%%%%%%%%%%%%%%%%%%%%%%%%%%%%%%%%%%%%%%%%%%%%%%

\subsection{Drought Monitoring System}

Drought is classified into several categories:

\begin{itemize}

\item Meteorological drought

\item Agricultural drought

\item Hydrological drought

\item Socioeconomic drought

\end{itemize}


The ACF Drought Engine integrates atmospheric, land and hydrological
information.

%%%%%%%%%%%%%%%%%%%%%%%%%%%%%%%%%%%%%%%%%%%%%%%%%%%%%%%%%%%%%%%%%%%%%%%%%%%%%%%%

\subsection{Standardized Precipitation Index (SPI)}

SPI measures precipitation anomalies over different time scales.

The index is:

\[
SPI=
\frac{P-\overline{P}}
{\sigma_P}
\]


where:

\begin{itemize}

\item $P$ = observed precipitation

\item $\overline{P}$ = climatological precipitation mean

\item $\sigma_P$ = precipitation standard deviation

\end{itemize}


SPI periods include:

\begin{itemize}

\item SPI-1 month

\item SPI-3 months

\item SPI-6 months

\item SPI-12 months

\end{itemize}


%%%%%%%%%%%%%%%%%%%%%%%%%%%%%%%%%%%%%%%%%%%%%%%%%%%%%%%%%%%%%%%%%%%%%%%%%%%%%%%%

\subsection{Standardized Precipitation Evapotranspiration Index}

SPEI includes the effect of temperature-driven evaporation.

\[
SPEI=
\frac{D-\overline{D}}
{\sigma_D}
\]


where:

\[
D=P-ET
\]


The system considers:

\begin{itemize}

\item Precipitation deficit

\item Temperature increase

\item Evapotranspiration demand

\end{itemize}


%%%%%%%%%%%%%%%%%%%%%%%%%%%%%%%%%%%%%%%%%%%%%%%%%%%%%%%%%%%%%%%%%%%%%%%%%%%%%%%%

\subsection{Palmer Drought Severity Index}

The Palmer Index evaluates long-term water balance.

It uses:

\begin{itemize}

\item Precipitation

\item Temperature

\item Soil moisture

\item Evapotranspiration

\end{itemize}


The ACF implementation combines classical drought indices with AI
corrections.

%%%%%%%%%%%%%%%%%%%%%%%%%%%%%%%%%%%%%%%%%%%%%%%%%%%%%%%%%%%%%%%%%%%%%%%%%%%%%%%%
%%%%%%%%%%%%%%%%%%%%%%%%%%%%%%%%%%%%%%%%%%%%%%%%%%%%%%%%%%%%%%%%%%%%%%%%%%%%%%%%
\subsection{Soil Moisture Monitoring}

Soil moisture is a critical variable linking atmospheric conditions,
land processes and drought development.

The ACF system monitors:

\begin{itemize}

\item Surface soil moisture

\item Root zone soil moisture

\item Soil water availability

\item Soil temperature

\item Freeze and thaw conditions

\end{itemize}


The soil water balance is represented by:

\[
\frac{dS}{dt}
=
P-E-R-D
\]

where:

\begin{itemize}

\item $S$ = soil water storage

\item $P$ = precipitation

\item $E$ = evapotranspiration

\item $R$ = runoff

\item $D$ = drainage

\end{itemize}


%%%%%%%%%%%%%%%%%%%%%%%%%%%%%%%%%%%%%%%%%%%%%%%%%%%%%%%%%%%%%%%%%%%%%%%%%%%%%%%%

\subsection{Vegetation Stress Monitoring}

Vegetation response is an important indicator of drought impact.

The ACF vegetation module analyzes:

\begin{itemize}

\item Leaf area index (LAI)

\item Vegetation health

\item Photosynthetic activity

\item Carbon exchange

\item Crop condition

\end{itemize}


Vegetation stress is estimated using:

\[
Stress=
f(Temperature,Water,Vegetation)
\]


The system detects:

\begin{itemize}

\item Crop water stress

\item Forest degradation

\item Ecosystem vulnerability

\item Agricultural drought

\end{itemize}


%%%%%%%%%%%%%%%%%%%%%%%%%%%%%%%%%%%%%%%%%%%%%%%%%%%%%%%%%%%%%%%%%%%%%%%%%%%%%%%%

\subsection{Satellite Drought Observation}

Satellite observations provide global drought monitoring capability.

The ACF Earth Observation System integrates:

\begin{itemize}

\item Land Surface Temperature

\item NDVI vegetation index

\item Soil moisture satellites

\item Evapotranspiration products

\item Surface water observations

\end{itemize}


Important satellite indicators:

\[
NDVI=
\frac{NIR-RED}
{NIR+RED}
\]


where:

\begin{itemize}

\item NIR = near infrared reflectance

\item RED = visible red reflectance

\end{itemize}


Low vegetation indices indicate possible vegetation stress.

%%%%%%%%%%%%%%%%%%%%%%%%%%%%%%%%%%%%%%%%%%%%%%%%%%%%%%%%%%%%%%%%%%%%%%%%%%%%%%%%

\subsection{Agricultural Drought Intelligence}

The ACF Agriculture Risk Module evaluates drought impacts on food
production.

The system analyzes:

\begin{itemize}

\item Crop development

\item Soil moisture deficit

\item Temperature stress

\item Rainfall anomalies

\item Growing season conditions

\end{itemize}


The agricultural risk index is:

\[
AgriculturalRisk=
f(SM,T,P,Vegetation)
\]


Outputs include:

\begin{itemize}

\item Crop stress maps

\item Yield risk estimation

\item Irrigation demand

\item Food security indicators

\end{itemize}


%%%%%%%%%%%%%%%%%%%%%%%%%%%%%%%%%%%%%%%%%%%%%%%%%%%%%%%%%%%%%%%%%%%%%%%%%%%%%%%%

\subsection{Artificial Intelligence Drought Prediction Engine}

The ACF AI Drought Engine improves early warning capability using
machine learning.

The hybrid approach is:

\[
Prediction=
ClimateModel+
Observation+
AI
\]


The AI system performs:

\begin{itemize}

\item Drought pattern recognition

\item Long-term anomaly detection

\item Seasonal forecasting

\item Impact prediction

\end{itemize}


%%%%%%%%%%%%%%%%%%%%%%%%%%%%%%%%%%%%%%%%%%%%%%%%%%%%%%%%%%%%%%%%%%%%%%%%%%%%%%%%

\subsection{Deep Learning Architecture}

The drought AI framework includes:

\begin{itemize}

\item Convolutional Neural Networks for climate maps

\item Transformers for temporal forecasting

\item Graph Neural Networks for Earth system connections

\item Physics-Informed Neural Networks

\end{itemize}


Input variables:

\begin{itemize}

\item Temperature fields

\item Precipitation

\item Soil moisture

\item Vegetation indices

\item Ocean climate signals

\end{itemize}


AI outputs:

\begin{itemize}

\item Drought probability

\item Severity classification

\item Forecast confidence

\item Future scenarios

\end{itemize}


%%%%%%%%%%%%%%%%%%%%%%%%%%%%%%%%%%%%%%%%%%%%%%%%%%%%%%%%%%%%%%%%%%%%%%%%%%%%%%%%

\subsection{Earth System Digital Twin Integration}

The Heatwave and Drought System is coupled with the ACF Earth System
Digital Twin.

The Digital Twin provides:

\begin{itemize}

\item Atmospheric state

\item Land surface conditions

\item Hydrological state

\item Climate projections

\item Ecosystem information

\end{itemize}


This coupling enables:

\begin{itemize}

\item Real-time drought monitoring

\item Heatwave evolution simulation

\item Climate change scenarios

\item Regional impact analysis

\end{itemize}


%%%%%%%%%%%%%%%%%%%%%%%%%%%%%%%%%%%%%%%%%%%%%%%%%%%%%%%%%%%%%%%%%%%%%%%%%%%%%%%%
%%%%%%%%%%%%%%%%%%%%%%%%%%%%%%%%%%%%%%%%%%%%%%%%%%%%%%%%%%%%%%%%%%%%%%%%%%%%%%%%
\subsection{Heatwave and Drought Risk Mapping}

The ACF Risk Mapping System transforms environmental indicators into
operational hazard information.

The risk calculation combines:

\begin{itemize}

\item Hazard intensity

\item Duration

\item Exposure

\item Vulnerability

\end{itemize}


The general risk formulation is:

\[
Risk =
Hazard
\times
Exposure
\times
Vulnerability
\]


The system generates:

\begin{itemize}

\item Global heatwave maps

\item Drought severity maps

\item Agricultural impact maps

\item Water stress maps

\item Human health risk maps

\end{itemize}


%%%%%%%%%%%%%%%%%%%%%%%%%%%%%%%%%%%%%%%%%%%%%%%%%%%%%%%%%%%%%%%%%%%%%%%%%%%%%%%%

\subsection{Human Health Impact Assessment}

Extreme heat creates direct risks for populations.

The ACF health impact module evaluates:

\begin{itemize}

\item Thermal stress

\item Population exposure

\item Urban heat islands

\item Vulnerable populations

\item Duration of extreme temperatures

\end{itemize}


The system provides:

\begin{itemize}

\item Heat warning levels

\item Exposure estimation

\item Emergency recommendations

\end{itemize}


%%%%%%%%%%%%%%%%%%%%%%%%%%%%%%%%%%%%%%%%%%%%%%%%%%%%%%%%%%%%%%%%%%%%%%%%%%%%%%%%

\subsection{Urban Heat Island Analysis}

Cities modify the natural thermal environment.

The ACF urban climate module analyzes:

\begin{itemize}

\item Building density

\item Surface materials

\item Vegetation coverage

\item Human activity

\item Night-time temperature

\end{itemize}


The urban heat anomaly is:

\[
UHI =
T_{urban}
-
T_{rural}
\]


where:

\begin{itemize}

\item $T_{urban}$ = urban temperature

\item $T_{rural}$ = surrounding rural temperature

\end{itemize}


%%%%%%%%%%%%%%%%%%%%%%%%%%%%%%%%%%%%%%%%%%%%%%%%%%%%%%%%%%%%%%%%%%%%%%%%%%%%%%%%

\subsection{AWCI Operational Integration}

The Atmospheric Weather Center Interface receives heatwave and drought
information from the ACF intelligence engine.

The dashboard displays:

\begin{itemize}

\item Temperature anomalies

\item Heatwave intensity

\item Drought indicators

\item Soil moisture maps

\item Vegetation stress

\item Agricultural risk

\item Warning levels

\end{itemize}


Operational workflow:

\[
Monitoring
\rightarrow
Detection
\rightarrow
Prediction
\rightarrow
RiskAssessment
\rightarrow
Warning
\]


%%%%%%%%%%%%%%%%%%%%%%%%%%%%%%%%%%%%%%%%%%%%%%%%%%%%%%%%%%%%%%%%%%%%%%%%%%%%%%%%

\subsection{Alert Management System}

The ACF Alert Engine generates dynamic warnings.

Alert categories:

\begin{itemize}

\item Heatwave warning

\item Extreme heat emergency

\item Agricultural drought alert

\item Water shortage alert

\item Ecosystem stress alert

\end{itemize}


The decision process considers:

\begin{itemize}

\item Forecast confidence

\item Population exposure

\item Duration

\item Expected impact

\end{itemize}


%%%%%%%%%%%%%%%%%%%%%%%%%%%%%%%%%%%%%%%%%%%%%%%%%%%%%%%%%%%%%%%%%%%%%%%%%%%%%%%%

\subsection{Forecast Verification}

The system continuously evaluates prediction performance.

Verification metrics include:

\begin{itemize}

\item Temperature forecast error

\item Drought index accuracy

\item Soil moisture prediction error

\item Vegetation response accuracy

\item Warning reliability

\end{itemize}


The forecast error is:

\[
Error =
Observed
-
Predicted
\]


%%%%%%%%%%%%%%%%%%%%%%%%%%%%%%%%%%%%%%%%%%%%%%%%%%%%%%%%%%%%%%%%%%%%%%%%%%%%%%%%

\subsection{Future Development}

Future ACF versions will include:

\begin{itemize}

\item Global kilometer-scale heatwave prediction

\item AI autonomous climate agents

\item Urban digital twins

\item Climate adaptation simulations

\item Integrated agriculture forecasting

\item Global water security prediction

\end{itemize}


%%%%%%%%%%%%%%%%%%%%%%%%%%%%%%%%%%%%%%%%%%%%%%%%%%%%%%%%%%%%%%%%%%%%%%%%%%%%%%%%

\subsection{References}

The scientific foundations include:

\begin{itemize}

\item Climate science

\item Land surface modeling

\item Hydrology

\item Remote sensing

\item Agricultural meteorology

\item Artificial Intelligence

\item Earth System Modeling

\end{itemize}


%%%%%%%%%%%%%%%%%%%%%%%%%%%%%%%%%%%%%%%%%%%%%%%%%%%%%%%%%%%%%%%%%%%%%%%%%%%%%%%%

